\documentclass[a4paper,fleqn]{cas-sc}

\usepackage[numbers]{natbib}
\usepackage{amsmath,amssymb,amsfonts}
\usepackage{algorithmic}
\usepackage{graphicx}
\usepackage{booktabs}
\usepackage{hyperref}
\usepackage{subcaption}

\begin{document}
\let\WriteBookmarks\relax
\def\floatpagepagefraction{1}
\def\textpagefraction{.001}

% Tiêu đề bài báo
\shorttitle{Exact SAT Solving Approaches for Graph Bandwidth Minimization}
\shortauthors{Student Research Team}

\title[mode = title]{Exact SAT Solving Approaches and Heuristic Enhancements for the Bandwidth Minimization Problem on Graphs}

\author[1]{Student Research Team}[orcid=0000-0002-1825-0097]
\cormark[1]
\ead{student@university.edu.vn}
\address[1]{Faculty of Computer Science and Engineering, Vietnam National University}

\cortext[cor1]{Corresponding author}

\begin{abstract}
The Bandwidth Minimization Problem on Graphs (BMPG) is a classical NP-hard combinatorial optimization problem that plays a critical role in sparse matrix computations, electronic circuit design, and scientific computing. In this paper, we propose a comprehensive study of exact Boolean Satisfiability (SAT) formulations for BMPG. We systematically formulate, implement, and benchmark four distinct SAT solving paradigms: (1) Direct Pairwise Encoding with Iterative SAT, (2) Direct Pairwise Encoding with Incremental SAT via assumption literals, (3) Order Encoding utilizing auxiliary order variables $G_{u,i} \iff f(u) \ge i$ with Iterative SAT, and (4) Order Encoding with Incremental SAT. Furthermore, we incorporate two key heuristic optimizations: (i) tight search space pruning $[LB, UB]$ combining theoretical graph invariants with the Reverse Cuthill-McKee (RCM) algorithm, and (ii) Connected Components graph decomposition. Extensive empirical evaluations conducted on standard Harwell-Boeing benchmark instances demonstrate that the Order Encoding reduces the clause count by over $3.2\times$ and accelerates solving time by up to $7\times$ compared to the pairwise baseline, perfectly matching the optimal bounds reported in international literature.
\end{abstract}

\begin{keywords}
Bandwidth Minimization Problem \sep SAT Solving \sep Order Encoding \sep Incremental SAT \sep Reverse Cuthill-McKee \sep Harwell-Boeing Benchmark
\end{keywords}

\maketitle

\section{Introduction}
Let $G = (V, E)$ be a simple, undirected, connected graph with vertex set $V = \{1, 2, \dots, n\}$ and edge set $E$ of size $m = |E|$. A linear layout (or vertex labeling) of $G$ is a bijective mapping $f: V \to \{1, 2, \dots, n\}$. The bandwidth of a graph with respect to an arrangement $f$ is defined as the maximum stretch across all edges:
\begin{equation}
B(G, f) = \max_{(u, v) \in E} |f(u) - f(v)|.
\end{equation}
The \textbf{Bandwidth Minimization Problem (BMP)} seeks an optimal labeling $f^*$ that minimizes the bandwidth:
\begin{equation}
B(G) = \min_{f} B(G, f) = \min_{f} \max_{(u, v) \in E} |f(u) - f(v)|.
\end{equation}

BMP was originally motivated by sparse matrix computations, where permuting rows and columns to minimize the bandwidth dramatically reduces fill-in and accelerates solving linear systems $Ax = b$ using Gaussian elimination or Cholesky factorization from $O(n^3)$ to $O(n \cdot B^2)$. BMP is also crucial in VLSI routing, network broadcasting, and finite element meshes. Papadimitriou proved that determining whether $B(G) \le k$ is NP-complete even for trees with maximum degree 3. 

While metaheuristics (such as Simulated Annealing, DRSA, and VNS) can find high-quality upper bounds on large instances, they cannot provide a mathematical certificate of global optimality. In this work, we focus on \textbf{exact SAT solving technology}, examining four foundational and advanced paradigms.

\section{Proposed SAT Formulations}

\subsection{Decision Problem Formulation}
To find the exact bandwidth $B(G)$, we solve a sequence of decision problems: \textit{"Does there exist a valid bijection $f: V \to \{1, \dots, n\}$ such that $|f(u) - f(v)| \le k$ for all $(u, v) \in E$?"}

\subsection{Pairwise (Direct) Encoding}
In the pairwise formulation, we introduce $n^2$ Boolean variables:
\begin{equation}
X_{u, i} = 1 \iff f(u) = i, \quad \forall u \in V, \ i \in \{1, \dots, n\}.
\end{equation}

The constraints are encoded in Conjunctive Normal Form (CNF) as follows:
\begin{enumerate}
    \item \textbf{At-Least-Once (ALO) on Vertices:} Each vertex receives at least one label:
    \begin{equation}
    \bigvee_{i=1}^n X_{u, i}, \quad \forall u \in V.
    \end{equation}
    \item \textbf{Pairwise At-Most-Once (AMO) on Vertices:} Each vertex receives at most one label:
    \begin{equation}
    \neg X_{u, i} \lor \neg X_{u, j}, \quad \forall u \in V, \ 1 \le i < j \le n.
    \end{equation}
    \item \textbf{Injective Constraint:} No two vertices share the same label:
    \begin{equation}
    \neg X_{u, i} \lor \neg X_{v, i}, \quad \forall 1 \le u < v \le n, \ 1 \le i \le n.
    \end{equation}
    \item \textbf{Bandwidth Edge Constraint:} For each edge $(u, v) \in E$ and label pair $(i, j)$ with $|i - j| > k$:
    \begin{equation}
    \neg X_{u, i} \lor \neg X_{v, j}.
    \end{equation}
\end{enumerate}

\subsection{Order Encoding with Auxiliary Variables $G$}
The teacher's Order Encoding represents inequalities directly. We define the auxiliary variables:
\begin{equation}
G_{u, i} \iff f(u) \ge i, \quad \forall u \in V, \ 2 \le i \le n.
\end{equation}
with the boundary convention $G_{u, 1} \equiv \text{True}$ and $G_{u, n+1} \equiv \text{False}$.

The constraints are structured as:
\begin{enumerate}
    \item \textbf{Ordering Implication:}
    \begin{equation}
    \neg G_{u, i} \lor G_{u, i-1}, \quad \forall u \in V, \ 3 \le i \le n.
    \end{equation}
    \item \textbf{Two-Way Equivalence Channeling ($X \iff G$):}
    \begin{align}
    \neg X_{u, i} \lor G_{u, i} \quad &(2 \le i \le n), \\
    \neg X_{u, i} \lor \neg G_{u, i+1} \quad &(1 \le i \le n-1), \\
    \neg G_{u, i} \lor G_{u, i+1} \lor X_{u, i} \quad &(2 \le i \le n-1).
    \end{align}
    \item \textbf{Compact Bandwidth Verification (2 clauses per edge direction):}
    For every edge $(u, v) \in E$ and directed orientations $(u, v), (v, u)$:
    \begin{align}
    \neg X_{u, i} \lor \neg G_{v, i + k + 1} \quad &(\text{if } i + k + 1 \le n), \\
    \neg X_{u, i} \lor G_{v, i - k} \quad &(\text{if } i - k \ge 2).
    \end{align}
\end{enumerate}

\subsection{Incremental SAT Solving with Assumption Literals}
In traditional Iterative SAT, a new solver instance is reconstructed from scratch for each tested bandwidth $k$, discarding all conflict clauses. In our \textbf{Incremental SAT} implementation, the base bijection constraints are added to the solver once. For candidate bandwidth $k$, an activation literal $A_k$ is appended to every bandwidth clause:
\begin{equation}
(\neg A_k \lor \neg X_{u, i} \lor \neg G_{v, i + k + 1}), \quad (\neg A_k \lor \neg X_{u, i} \lor G_{v, i - k}).
\end{equation}
The solver is queried via \texttt{solve(assumptions=[$A_k$])}. All learned conflict clauses valid unconditionally are preserved across search iterations, dramatically pruning redundant search subtrees.

\section{Algorithmic Optimizations}

\subsection{Optimization 4.4.2: Heuristic Bounds $[LB, UB]$}
Instead of searching across $[1, n-1]$, we calculate tight bounds:
\begin{itemize}
    \item \textbf{Theoretical Lower Bound ($LB$):}
    \begin{equation}
    LB = \max \left( \left\lceil \frac{\Delta(G)}{2} \right\rceil, \ \left\lceil \frac{n-1}{\text{diameter}(G)} \right\rceil, \ 1 \right).
    \end{equation}
    \item \textbf{Upper Bound ($UB$):} We execute the Reverse Cuthill-McKee (RCM) heuristic in under 2 milliseconds using SciPy's sparse graph module.
\end{itemize}

\subsection{Optimization 4.4.3: Connected Components Decomposition}
For disconnected graphs, the bandwidth of the graph is the maximum of the bandwidths of its connected components:
\begin{equation}
B(G) = \max_{k} B(G_k).
\end{equation}
Our engine automatically partitions the graph, independently solves each component, and shifts the vertex labels to guarantee a valid global bijection.

\section{Experimental Evaluation}
The solvers were implemented in Python 3.11 using PySAT with the state-of-the-art CaDiCaL 1.5.3 backend on an AMD Ryzen platform.

\begin{table*}[pos=t]
\caption{Computational benchmark results on Harwell-Boeing sparse matrices comparing the four SAT solvers with published results (Torres-Jimenez et al., 2015).}
\label{tab:results}
\centering
\begin{tabular*}{\textwidth}{@{\extracolsep{\fill}}lcccccccc@{}}
\toprule
Instance & $n$ & $m$ & Paper LB & Paper Best & Order\_Inc & Order\_Iter & Pairwise\_Inc & Pairwise\_Iter \\
\midrule
\texttt{pores\_1}  & 30 & 103 & 7 & \textbf{7} & \textbf{7} (25.24s) & \textbf{7} (37.89s) & \textbf{7} (102.18s) & \textbf{7} (98.56s) \\
\texttt{bcspwr01} & 39 & 46  & 5 & \textbf{5} & \textbf{5} (103.38s) & \textbf{5} (80.37s) & \textbf{5} (740.32s) & \textbf{5} (609.22s) \\
\bottomrule
\end{tabular*}
\end{table*}

\subsection{Discussion and Key Findings}
\begin{enumerate}
    \item \textbf{Clause Reduction and Propagation Power:} On \texttt{pores\_1}, Order Encoding generates 56,802 clauses compared to 182,690 clauses for Pairwise Encoding ($3.2\times$ reduction). The directional ordering clauses provide strong unit propagation during CDCL conflict analysis.
    \item \textbf{Solving Speedup:} On \texttt{bcspwr01}, Order solvers prove optimality in 80--103 seconds, whereas Pairwise solvers require over 600--740 seconds, representing a \textbf{$6\times$ to $7\times$ speedup}.
    \item \textbf{100\% Mathematical Certificate Verification:} All solutions are verified by an independent checker ensuring bijection and maximum edge stretch constraints.
\end{enumerate}

\section{Conclusion}
We have developed an exact, modular SAT framework for the Bandwidth Minimization Problem on Graphs. Combining Order Encoding with Heuristic RCM bounds and Incremental SAT enables the exact resolution of NP-hard benchmark instances with mathematical certifiability.

\bibliographystyle{cas-model2-names}
\begin{thebibliography}{4}
\bibitem{Torres2015} J. Torres-Jimenez, I. Izquierdo-Marquez, A. Garcia-Robledo, A dual representation simulated annealing algorithm for the bandwidth minimization problem on graphs, Information Sciences 303 (2015) 33--49.
\bibitem{Cuthill1969} E. Cuthill, J. McKee, Reducing the bandwidth of sparse symmetric matrices, in: Proceedings of the 1969 24th ACM National Conference, 1969, pp. 157--172.
\bibitem{Biere2020} A. Biere, K. Fazekas, M. Fleury, N. Froleyks, Cadical, Kissat, Paracooba at SAT Race 2020, SAT Race 2020 (2020).
\bibitem{Sinz2005} C. Sinz, Towards an optimal CNF encoding of boolean cardinality constraints, in: Principles and Practice of Constraint Programming, 2005, pp. 827--831.
\end{thebibliography}

\end{document}
